\documentclass{article}
\usepackage[T1]{fontenc}
\usepackage{lmodern}
\usepackage{graphicx}
\usepackage{amsmath,amssymb}
\usepackage[a4paper,margin=2cm]{geometry}
\usepackage[absolute,overlay]{textpos}
\setlength{\TPHorizModule}{1mm}
\setlength{\TPVertModule}{1mm}
\usepackage[indonesian]{babel}
\usepackage{hyperref}
\setlength{\emergencystretch}{2em}
% unit: Halaman 18

\begin{document}

% === Halaman 18 (python/native) ===

Blok plainteks ke-1 

: 011000 

 Kunci publik  

: 62, 93, 81, 88, 102, 37

 Kriptogram 


: (1  93) + (1  81)  = 174

Blok plainteks ke-2 

: 110101 

 Kunci publik  

: 62, 93, 81, 88, 102, 37

 Kriptogram 


: (1  62) + (1  93) + (1  88) + 




  (1  37)  = 280

Blok plainteks ke-3 

: 101110 

 Kunci publik  

: 62, 93, 81, 88, 102, 37

 Kriptogram 


: (1  62) + (1  81) + (1  88) + 




  (1  102)  = 333

 Jadi, cipherteks yang dihasilkan : 174, 280, 333

18

\end{document}
